% 原书分章 PDF 第 21 页：第 8.8 节。
\section{讨论}

\begin{sourcebox}
区间删失数据与双变量右删失数据的例子，代表了有效影响曲线不存在闭式表达的一类问题；相应地，构造有效估计量一直极具挑战。这些例子说明，可以用 HAL 最小二乘线性回归估计这些复杂的有效影响曲线，使一步估计量或 TMLE 的构造变得相对直接。这是一个相当出色的结果！
\end{sourcebox}

在带协变量的双变量生存函数例子中，我们说明：对给定 $(A,W)$ 使用这种 Lasso 回归估计量，甚至可以估计在控制高维协变量向量 $W$ 后，二值处理 $A$ 对双变量生存函数的因果效应的有效影响曲线。这样，就有可能根据双变量右删失数据，为这一双变量生存函数上的因果效应开发 TMLE。对于扩展区间删失数据结构
$O=(W,A,\Delta_1,C_1,\ldots,\Delta_M,C_M)$，
我们还用相同方法，展示了二值处理对生存时间截断均值的因果效应的有效影响曲线。这使我们能够估计处理对受到区间删失的事件发生时间的因果效应。同样，后一问题也是一个很有意义、具有重要实际应用价值的估计问题，尚需未来继续研究。我们的 $D^*(P)$ 的 HAL 估计量依赖来自 $P$ 的独立同分布样本。不过，我们还说明（未给出正式证明），如果目标是估计 $D^*(P_0)$，而 $P$ 表示 $P_0$ 的一致估计量，那么，也可以直接对来自 $P_0$ 的原始样本应用 HAL 估计量。这使所提出的 HAL 估计量格外方便。

本章提出的有效影响曲线估计方法，为 Frangakis 等（2015）、Luedtke 等（2015a）以及 van der Laan 等（2015）提出的方法提供了另一种选择。后一类研究由 Frangakis 等（2015）富有启发性的文章开启，旨在使未受过有效性理论训练的人，也能通过计算机估计有效影响曲线，进而进行有效估计。显然，本章的方法仍要求使用者提出初始梯度、得分算子和相应回归问题。不过，它避免了寻找典范梯度的闭式表达，也无需从数学上或数值上计算初始梯度在切空间上的投影，而是利用先进的预测机器学习方法来近似这一投影。
